\chapter{逆变电路}

把直流电变成交流电称为\textbf{逆变}。逆变电路与整流电路相对应，是电力电子技术中最重要的变换之一：整流是把交流“整理”成直流，逆变则是把直流“还原”成交流。按交流侧的去向，逆变分为两类——\textbf{有源逆变}（交流侧接电网，已见 3.7 节）与\textbf{无源逆变}（交流侧直接接负载，如变频器带电动机、不间断电源带负载）。本章讲的是\textbf{无源逆变}。

逆变电路的核心问题是\textbf{换流}：电流要从一个桥臂转移到另一个桥臂，能不能可靠转移，决定了电路能不能工作。

% ============================================================
\section{换流方式}

\textbf{换流}：电流从一个桥臂向另一个桥臂转移的过程。按促使晶闸管（或全控器件）关断的机理不同，换流方式分为四种。

\subsection{器件换流}

利用\textbf{全控型器件}的自关断能力进行换流。由于 IGBT、电力 MOSFET、GTO、GTR 等器件都能由控制信号关断，因此由它们构成的电路全部采用器件换流。

\begin{itemize}
  \item 优点：不需要附加换流电路，主电路简单、效率高；
  \item 这是现代逆变电路的主流方式（第 4 章后面各节，除特别说明外均为器件换流）。
\end{itemize}

\subsection{电网换流}

由\textbf{电网}提供换流电压——只要把负的电网电压加在欲关断的晶闸管上即可使其关断。这种方式不需要器件具有自关断能力，也不需要附加换流电路，经济方便；但\textbf{只能用于由交流电网供电的电路}（如相控整流电路、有源逆变电路、交交变频电路），不能用于无源逆变。

\subsection{负载换流}

由\textbf{负载}提供换流电压。条件是\textbf{负载电流超前于负载电压}（容性负载），这样在电流过零之前电压已经反向，使欲关断的晶闸管承受反压而关断。

典型应用是\textbf{并联谐振式逆变电路}（用于感应加热、中频电源）：负载为感应线圈（$L$、$r$ 串联）与补偿电容 $C$ 并联，电路工作在接近谐振的状态，负载对基波呈高阻抗、对谐波呈低阻抗，因此负载电压接近正弦波。为保证可靠换流，触发脉冲必须\textbf{引前}于负载电压一个角度 $\beta$（逆变角／触发引前角），

\keyformula{\beta>\gamma+\delta}

式中 $\gamma$ 为换相重叠角，$\delta$ 为晶闸管关断时间所对应的电角度。$\beta$ 太小则换流失败（逆变颠覆）；$\beta$ 太大则输出功率因数下降、输出电压降低。实际取 $\beta$ 略大于 $\gamma+\delta$ 的最小值。

\subsection{强迫换流}

设置\textbf{附加的换流电路}，由换流电路内的电容（或电感）的储能向欲关断的晶闸管施加反向电压或反向电流使其关断。由于换流电压／电流是“强行”加上的，故称强迫换流，也称脉冲换流。

\begin{itemize}
  \item 直接耦合式强迫换流：由换流电路内电容直接提供换流电压；
  \item 电感耦合式强迫换流：通过换流电路内电容和电感耦合提供换流电压或换流电流。
\end{itemize}

强迫换流需要附加电路，增加了成本与损耗，但对晶闸管电路而言是唯一可用于无源逆变的换流方式（晶闸管本身不能自关断，又没有电网或超前负载电压可用）。

\begin{table}[htbp]
\centering
\caption{四种换流方式的比较}
\label{tab:4-1}
\begin{tabular}{p{2.4cm}p{4.4cm}p{7.6cm}}
\toprule
\textbf{换流方式} & \textbf{换流能量来源} & \textbf{适用场合} \\
\midrule
器件换流 & 全控器件的控制信号 & 由 IGBT、MOSFET、GTO、GTR 构成的各类电路 \\
电网换流 & 交流电网电压 & 相控整流、有源逆变、交交变频（必须接电网） \\
负载换流 & 负载的超前电压 & 并联谐振式逆变、感应加热中频电源（负载须容性） \\
强迫换流 & 附加电路的储能（电容/电感） & 晶闸管无源逆变电路（需附加换流电路） \\
\bottomrule
\end{tabular}
\end{table}

% ============================================================
\section{电压型逆变电路}

\textbf{电压型逆变电路}（Voltage Source Inverter，VSI）的直流侧是\textbf{电压源}，用大电容储能并使直流电压基本无脉动，因此交流侧输出电压为\textbf{矩形波}（或阶梯波），而电流波形取决于负载。

\begin{itemize}
  \item 直流侧并联大电容，相当于电压源，直流侧电压基本无脉动，直流回路呈现\textbf{低阻抗}；
  \item 交流侧输出电压波形为矩形波，与负载阻抗角无关；输出电流波形和相位因负载不同而不同；
  \item 需要器件换流（全控器件）或强迫换流（晶闸管）；
  \item \textbf{当交流侧为阻感负载时需要提供无功功率}，直流侧电容起缓冲无功能量的作用；为给交流侧向直流侧反馈的无功能量提供通道，各桥臂的开关器件\textbf{都须反并联二极管（续流二极管）}。
\end{itemize}

\subsection{单相电压型逆变电路}

\subsubsection*{1. 半桥逆变电路}

由两只开关器件（$S_1$、$S_2$）与两只大电容串联组成，负载接在两电容中点和两器件中点之间。$S_1$、$S_2$ 交替导通、各导通 $180^\circ$，输出方波电压的幅值为 $U_d/2$。

\keyformula{U_o=\frac{U_d}{2},\qquad U_{o1m}=\frac{2U_d}{\pi},\qquad U_{o1}=\frac{\sqrt{2}U_d}{\pi}=0.45U_d}

其中 $U_{o1m}$ 为基波幅值、$U_{o1}$ 为基波有效值。半桥电路使用器件少，但直流侧需两只电容分压，且实际中两电容电压可能不均衡（需均压措施），故只用于几千瓦以下的小功率场合。

\subsubsection*{2. 全桥逆变电路}

由四个开关器件组成两个桥臂，可看成两个半桥电路的组合。$S_1$、$S_4$ 为一对（同时导通），$S_2$、$S_3$ 为另一对，两对交替各导通 $180^\circ$，输出方波电压幅值为 $U_d$（是半桥的两倍）。

\keyformula{U_o=U_d,\qquad U_{o1m}=\frac{4U_d}{\pi},\qquad U_{o1}=\frac{2\sqrt{2}U_d}{\pi}=0.9U_d}

\begin{figure}[htbp]
\centering
\begin{tikzpicture}[font=\footnotesize,scale=0.95]
  \draw[->] (0,0) -- (8.4,0) node[below] {$\omega t$};
  \draw[->] (0,-2.2) -- (0,2.6) node[left] {$u_o$};
  \draw[yhbAccent,thick] (0,2.0) -- (4,2.0) -- (4,-2.0) -- (8,-2.0) -- (8,2.0) -- (8.3,2.0);
  \draw[dashed,textgray] (0,2.0) node[left] {\scriptsize $U_d$} -- (0,2.0);
  \draw[dashed,textgray] (0,-2.0) node[left] {\scriptsize $-U_d$};
  \draw[yhbAccent2,thick] plot[domain=0:8,samples=120] (\x,{2*4/3.1416*sin(57.3*\x)});
  \node[text=yhbAccent2] at (6.8,2.5) {\scriptsize 基波};
  \node[text=yhbAccent] at (0.75,2.35) {\scriptsize 方波};
\end{tikzpicture}
\caption{单相全桥电压型逆变电路的输出电压波形}
\label{fig:4-1}
\end{figure}

方波电压含有的谐波为 $n=2k\pm1$（奇次），各次谐波幅值与次数成反比：$U_{onm}=\dfrac{4U_d}{n\pi}$。因此\textbf{要改变输出电压有效值，只能改变直流电压 $U_d$}；要减少谐波，则须采用 PWM 技术（第 7 章）或多重化。

\subsection{三相电压型逆变电路}

三个桥臂、六只开关器件组成三相桥，是应用最广的逆变电路。常用\textbf{180$^\circ$ 导电方式}：每个桥臂的上下两管互补导通各 $180^\circ$，三个桥臂的导通依次相差 $120^\circ$，因此\textbf{任何时刻都有三个桥臂同时导通}（同一桥臂上下管绝不能同时导通，否则直流侧短路，实际须留“死区”）。

输出波形：

\begin{itemize}
  \item \textbf{线电压}为 $120^\circ$ 矩形波（幅值 $U_d$），因为线电压只取决于两相桥臂的状态；
  \item \textbf{相电压}（星形负载）为\textbf{六阶梯波}，幅值为 $\pm2U_d/3$ 与 $\pm U_d/3$ 交替。
\end{itemize}

\keyformula{U_{UN}=\frac{\sqrt{2}}{3}U_d=0.471U_d,\quad U_{UN1m}=\frac{2U_d}{\pi}=0.637U_d,\quad U_{UN1}=\frac{\sqrt{2}U_d}{\pi}=0.45U_d}

\keyformula{U_{UV}=\sqrt{\frac{2}{3}}U_d=0.816U_d,\quad U_{UV1m}=\frac{2\sqrt{3}U_d}{\pi}=1.1U_d,\quad U_{UV1}=\frac{\sqrt{6}U_d}{\pi}=0.78U_d}

\corepoint{三相电压型桥式逆变的三个“0.471 / 0.816 / 0.45 / 0.78”必须记准：\textbf{相电压有效值 0.471$U_d$、线电压有效值 0.816$U_d$、相电压基波有效值 0.45$U_d$、线电压基波有效值 0.78$U_d$}。还要记住两点：（1） 同一桥臂上下管\textbf{互补导通}，换流是“器件换流”，但必须设\textbf{死区}防止直通；（2） 负载为阻感时，无功能量经\textbf{反并联二极管}回馈到直流侧电容。}

% ============================================================
\section{电流型逆变电路}

\textbf{电流型逆变电路}（Current Source Inverter，CSI）的直流侧是\textbf{电流源}，用大电感储能并使直流电流基本无脉动，因此交流侧输出电流为\textbf{矩形波}，而电压波形取决于负载。

\begin{itemize}
  \item 直流侧串联大电感，相当于电流源，直流回路呈现\textbf{高阻抗}；
  \item 交流侧输出电流波形为矩形波，与负载阻抗角无关；输出电压波形和相位因负载不同而不同；
  \item 因为直流侧电感储能，\textbf{无功能量由电感缓冲}，不需要反并联二极管（这一点与电压型相反）；
  \item 换流时须采取\textbf{“先断后通”}的方法：先给应关断的器件关断信号，待其关断并留一定时间裕量后，再给应导通的器件发开通信号，否则将造成直流侧电源短路。
\end{itemize}

\subsection{单相电流型逆变电路}

典型电路是\textbf{并联谐振式逆变电路}（用于感应加热、中频电源）：四个桥臂的晶闸管各串联一个电抗器 $L_T$（限制开通时的 di/dt），负载为感应线圈（$L$、$r$ 串联）与补偿电容 $C$ 并联，构成并联谐振电路。

\begin{itemize}
  \item 工作频率接近负载电路的谐振频率，负载对基波呈\textbf{高阻抗}、对谐波呈低阻抗，故\textbf{负载电压接近正弦波}；
  \item 换流方式为\textbf{负载换流}：负载电流超前于负载电压（容性），电流过零前电压已反向，使欲关断的晶闸管承受反压而关断，因此必须使触发脉冲\textbf{引前}负载电压一个角度 $\beta$；
  \item 由于 $L_T$ 的存在，换相需要时间，形成\textbf{换相重叠角} $\gamma$；负载的功率因数角 $\varphi$ 与 $\beta$、$\gamma$ 的关系决定输出电压大小，$\beta$ 越大输出电压越低（功率因数越低）；
  \item 起动问题：需先给电容充电（预充电），起动时把电容能量释放到负载形成衰减振荡，检测出振荡信号实现自励。
\end{itemize}

\subsection{三相电流型逆变电路}

典型电路是三相桥式电流型逆变电路，采用\textbf{120$^\circ$ 导电方式}：每个桥臂导通 $120^\circ$，任何时刻只有\textbf{两个桥臂同时导通}（一个共阴组、一个共阳组），输出电流为 $120^\circ$ 矩形波。

\keyformula{I=\sqrt{\frac{2}{3}}I_d=0.816I_d,\qquad I_1=\frac{\sqrt{6}I_d}{\pi}=0.78I_d}

其中 $I$ 为输出电流有效值，$I_1$ 为基波有效值。各桥臂之间换流采用\textbf{强迫换流}，连接于各桥臂之间的电容 $C_1\sim C_6$ 即为换流电容；换流过程分“恒流放电阶段”（放电电流恒为 $I_d$）与“谐振放电阶段”等。由于是强迫换流，换流电容的容量必须足以关断晶闸管，故工作频率不能太高。

\begin{table}[htbp]
\centering
\caption{电压型与电流型逆变电路的比较}
\label{tab:4-2}
\begin{tabular}{p{2.6cm}p{6.0cm}p{6.0cm}}
\toprule
\textbf{比较项} & \textbf{电压型（VSI）} & \textbf{电流型（CSI）} \\
\midrule
直流侧储能元件 & 大电容（并联），电压源性质、低阻抗 & 大电感（串联），电流源性质、高阻抗 \\
输出波形 & 输出电压为矩形波，电流取决于负载 & 输出电流为矩形波，电压取决于负载 \\
反并联二极管 & 必须有（为无功能量提供通路） & 不需要 \\
换流方式 & 器件换流（或强迫换流） & 负载换流、强迫换流；须“先断后通” \\
短路保护 & 直流侧短路电流大，保护较难 & 电感限流，短路电流上升慢，保护容易 \\
动态响应 & 快 & 较慢（电感电流不能突变） \\
典型应用 & 通用变频器、UPS、开关电源 & 感应加热、中频电源、大容量电动机调速 \\
\bottomrule
\end{tabular}
\end{table}

% ============================================================
\section{多重逆变电路和多电平逆变电路}

前面介绍的逆变电路输出都是矩形波，含大量低次谐波。改善输出波形的途径有两条：一是\textbf{把几个矩形波按一定相位差组合起来}，使某些主要谐波相互抵消——\textbf{多重逆变电路}；二是\textbf{改变电路结构，输出更多的电平}，使波形向正弦波靠近——\textbf{多电平逆变电路}。

\subsection{多重逆变电路}

把若干个逆变电路的输出按一定的相位差组合起来，使它们所含的某些主要谐波分量相互抵消，就可以得到较为接近正弦波的波形。

\begin{itemize}
  \item \textbf{二重逆变电路}：两个单相全桥逆变电路（输出电压相位错开一定角度）经变压器合成，输出波形为阶梯波，最低次谐波次数提高、幅值减小；两个三相桥式逆变电路相位错开 $30^\circ$ 合成，可构成 \textbf{12 脉波}逆变电路。
  \item 一般规律：使 $m$ 个三相桥式逆变电路的相位依次错开 $\pi/(3m)$ 运行，连同使它们输出电压合成并抵消上述相位差的变压器，就可以构成\textbf{脉波数为 $6m$ 的逆变电路}（二重 12 脉波、三重 18 脉波）。
  \item 多重化的效果：输出波形中最低次谐波次数提高到 $6m\pm1$，谐波幅值随次数成反比减小，故波形明显改善；同时每个逆变单元承受的功率减小，便于用较低电压等级的器件构成大容量装置。
\end{itemize}

\subsection{多电平逆变电路}

\textbf{多电平逆变电路}能够输出较多的电平，从而使输出电压向正弦波靠近。常见结构有三种：

\begin{enumerate}
  \item \textbf{中点钳位型（NPC）三电平逆变电路}：每相桥臂由四个开关器件串联，两只钳位二极管把桥臂中点与直流侧电容中点相连。以 U 相为例，当 $V_{11}$、$V_{12}$（或 $VD_{11}$、$VD_{12}$）导通、$V_{41}$、$V_{42}$ 关断时，U 点与直流中点（$\varphi$ 点）之间的电压为 $+U_d/2$；反之则为 $-U_d/2$；上下各一只导通时输出为 0。因此每相可输出 $+U_d/2$、$0$、$-U_d/2$ 三种电平，线电压有五种电平。
  \item \textbf{飞跨电容型（FLC）三电平逆变电路}：用飞跨电容代替钳位二极管，结构稍简单，但需要控制电容电压平衡，使用较少；构成更多电平时所需电容数目急剧增加。
  \item \textbf{单元串联多电平逆变电路}：把多个独立的低压逆变单元（H 桥）串联叠加，用移相变压器分别供电，可获得很高的输出电压与很多电平数，谐波小、可靠性高（单元故障可旁路），在高压大容量变频器中应用较多。
\end{enumerate}

\corepoint{本章的考点集中在三处：\textbf{（1） 四种换流方式的判别}（看到“全控器件”想到器件换流，看到“接电网”想到电网换流，看到“负载电流超前”想到负载换流，看到“附加电容”想到强迫换流）；\textbf{（2） 电压型与电流型的对比}（电容／电感、是否要反并联二极管、是否“先断后通”）；\textbf{（3） 输出电压有效值与基波的四个系数}（单相全桥 $0.9U_d$、三相相电压 $0.45U_d$ 与 $0.471U_d$、线电压 $0.78U_d$ 与 $0.816U_d$）。要改输出电压有效值只能改 $U_d$，要改波形只能靠 PWM 或多重化。}
